% ------------------------------------------------------------------------------------
%                       PROGRAMMAZIONE LINEARE INTERA
% ------------------------------------------------------------------------------------
\chapter{Problemi di modellazione}

\section{Knapsack}

Il problema è il seguente:

\begin{exa}[KNAPSACK]
    Immaginiamo di avere
    \begin{itemize}
        \item $n$ \textbf{oggetti}, ciascuno con
            \begin{itemize}
                \item Un \textbf{peso} $w_j$.
                \item Un \textbf{profitto} $p_j$
            \end{itemize}
            Con $j\in J = \{ 1,\ldots ,n \}  $
        \item Un \textbf{contenitore} (\textit{knapsack}) 
            di \textbf{capacita massima} $W$.
    \end{itemize}
    Vogliamo capire come \textbf{massimizzare il profitto} 
    rimanendo \textbf{entro i limiti di capacita}.
\end{exa}

Introduciamo le \textbf{variabili binarie} seguenti:
\[
    x_j = \begin{cases}
        1 & \text{se item j viene messo nel contenitore} \\ 
        0 & \text{altrimenti}
    \end{cases}
\]

Il problema puo quindi essere formalizzato come segue:
\[
    \begin{cases}
        \text{max } \sum_{j \in J} p_{j} x_{j} \\ 
        \sum_{j\in J} w_{j} x_{j} \le W \\ 
        x_{j} \in \{0,1\} & \forall j \in J
    \end{cases}
\]

\section{Facility location}

Il problema è il seguente:

\begin{exa}[FACILITY LOCATION]
    Abbiamo:
    \begin{itemize}
        \item Un \textbf{insieme} $J$ di $n$ \textbf{utenti}.
        \item Un \textbf{insieme} $I$ di $m$ \textbf{facilities}, ciascuna
            con un \textbf{costo fisso di costruzione} $d_i$.
        \item Ciascun \textbf{utente j comporta} un \textbf{costo di servizio}
            per poter \textbf{usufruire} della \textbf{i-esima facility} 
            pari a $c_{ij}$.
    \end{itemize}
    Vogliamo \textbf{pianificare} la \textbf{costruzione delle facilities}
    in modo da \textbf{minizzare i costi di servizio}.
\end{exa}

Introduciamo le variabili binarie seguenti:
\[
    y_i = \begin{cases}
        1 & \text{se costruiamo la facility i} \\ 
        0 & \text{altrimenti}
    \end{cases}
\]

Ed anche:
\[
    x_{ij} = \begin{cases}
        1 & \text{se utente j e servito dalla facility i} \\ 
        0 & \text{altrimenti}
    \end{cases}
\]

Il modello risulta pertanto essere:
\[
    \begin{cases}
        \text{min } \sum_{i\in I} d_i y_i + 
        \sum_{i \in I} \sum_{j \in J} c_{ij} x_{ij} \\ 
        \sum_{i \in I} x_{ij} = 1 \\ 
        x_{ij} \le y_i & \forall i \in I \text{ , } \forall j \in J \\
        y_i \in \{0,1\} & \forall i \in I \\ 
        x_{ij} \in \{0,1\} & \forall i \in I \text{ , } \forall j \in J  
    \end{cases}
\]

\begin{remark}[VINCOLI]
    Il vincolo
    \[
        x_{ij} \le y_i \text{ } \forall i \in I \text{ , } \forall j \in J
    \] 
    sta ad indicare il fatto che:
    \[
        \begin{aligned}
            y_i = 0 &\implies x_{ij} = 0 \text{ , } \forall j \in J \\ 
            x_{ij} = 1 &\implies y_i = 1 \text{ , } \forall i \in I
        \end{aligned}
    \] 
    In \textbf{totale}, questi vincoli sono $n \cdot m$ in \textbf{numero}.
\end{remark}

Ci chiediamo allora se \textit{possiamo fare di meglio}, ovvero
se possiamo \textit{rappresentare la stessa implicazione}
usando \textit{meno vincoli}.
La risposta è \textbf{sni}: possiamo scrivere i vincoli come
\[
    \sum_{j\in J} x_{ij} \le ny_i \text{ , } \forall i \in I
\]

che però comporta un \textbf{problema pratico}:

\begin{remark}[COMPROMESSO]
    Abbiamo \textbf{ridotto} il \textbf{numero di questi vincoli} 
    da $n \cdot m$ a $m$.

    \textbf{Tuttavia}, questa soluzione \textbf{risulta peggiore in pratica}
    in quanto \textbf{peggiora l'efficienza dei rilassamenti}. 
    Inoltre, e facile intuire che i \textbf{vincoli aggregati} che
    abbiamo scritto sarebbero stati \textbf{deducibili anche dai vincoli singoli}.
\end{remark}

\section{Set covering}

Il problema è il seguente:

\begin{exa}[SET COVERING]
    Abbiamo come \textbf{input}:
    \begin{itemize}
        \item un \textbf{insieme} $I$ detto \textbf{ground set}.
        \item Una \textbf{famiglia} $\mathcal{F} = \{ F_{1}, \ldots , F_n\}$ 
            dove $F_j \subseteq I$ $\forall j \in J = \{1,\ldots ,n\}  $.
    \end{itemize}
    Vogliamo \textbf{scegliere un sottoinsieme} $\mathcal{F}^* = \{F_k\}$ 
    con $k \in K \subseteq J$ di $\mathcal{F}$ tale per cui
    \[
        \bigcup_{k \in K} F_k = I
    \] 
\end{exa}

Definiamo le \textbf{variabili binarie} seguenti:
\[
    x_j = \begin{cases}
        1 & \text{se scelgo } F_j \\ 
        0 & \text{altrimenti}
    \end{cases}
\]

Il modello risulta quindi essere:
\[
    \begin{cases}
        \text{min } \sum_{j \in J} x_j \\ 
        \sum_{j \text{ : } i \in F_j} x_j \ge 1 & \forall i \in I \\ 
        x_j \in \{0,1\} & \forall j \in J  
    \end{cases}
\]

\begin{remark}[APPLICAZIONI REALI]
    Problemi di questo tipo sono \textbf{risolti giornalmente}
    dalle \textbf{compagnie aeree} per \textbf{determinare efficientemente}
    i \textbf{turni} dei \textbf{piloti}.
\end{remark}

\begin{remark}[SET PARTIONING]
    A volte il problema puo essere modellato come
    problema di \textbf{set partionining}, in cui il
    \textbf{ground set} deve esere \textbf{coperto senza sovrapposizioni}.
    Chiaramente, questa formulazione risulta \textit{piu difficile da risolvere}.
\end{remark}
